\documentclass[10pt,letterpaper]{article}
\usepackage{cornell-notes}
\usepackage{mathtools}

\title{Chapter 18: Two-Point Boundary Value Problems}
\author{}
\date{}

\setDocTitle{Chapter 18: Two-Point Boundary Value Problems}
\setDocSubtitle{Numerical Methods}
\setDocVersion{v1.0}
\setDocStatus{Study Companion}
\setDocOwner{Numerical Methods Cornell Notes}
\setDocAuthor{}
\setDocDate{\today}
\setDocSummary{Cornell-format study and review notes for Chapter 18: Two-Point Boundary Value Problems.}

\setCornellCollection{Numerical Methods Cornell Notes}
\setCornellUnitType{Chapter}
\setCornellUnitNumber{18}
\setCornellUnitTitle{Two-Point Boundary Value Problems}
\setCornellCompanionLabel{Study and review companion}

\hypersetup{pdftitle={Chapter 18: Two-Point Boundary Value Problems},pdfsubject={Cornell notes},pdfauthor={}}

\begin{document}
\maketitle

\section*{Chapter Overview}
This chapter organizes the mathematical and computational ideas behind two-point boundary value problems. The notes preserve the numbered section sequence while emphasizing method selection, explicit assumptions, numerical reliability, cost, and verification.

\section*{Learning Objectives}
Explain the governing problem class; compare The Shooting Method, Shooting to a Fitting Point, Relaxation Methods, and A Worked Example: Spheroidal Harmonics; design defensible stopping and verification checks; distinguish problem conditioning from algorithmic stability.

\section*{Prerequisites}
ODE integration, nonlinear systems, finite differences, and boundary conditions.

\section*{Operational Workflow}
Count differential order and independent boundary constraints; choose shooting, multiple shooting, matching, or relaxation; provide initial guesses; solve endpoint or mesh residuals; refine and verify conditions.

\subsection*{Formula and notation anchor}
\[
y'(x)=f(x,y),\qquad B\bigl(y(a),y(b)\bigr)=0
\]
Unlike an initial-value problem, the boundary conditions couple unknown data at separated locations and generally require an outer solve.

\begin{warnbox}{Numerical warning}
Exponentially unstable shooting, singular endpoints, inconsistent boundary counts, poor initial guesses, mesh underresolution, ill-conditioned Jacobians, and arbitrary normalization of eigenfunctions.
\end{warnbox}

\section{18.0 Introduction}
\begin{cornell}
\noterow{What is the central idea?}{A boundary-value problem imposes conditions at separated points, so the entire solution trajectory must satisfy coupled endpoint constraints.}
\noterow{How should it be used?}{For Introduction, begin from explicit inputs, outputs, preconditions, and representation choices.  Count differential order and independent boundary constraints; choose shooting, multiple shooting, matching, or relaxation; provide initial guesses; solve endpoint or mesh residuals; refine and verify conditions.}
\noterow{Cost and reliability}{Analyze arithmetic work, storage, data movement, and convergence separately.  Relevant failure pressures include exponentially unstable shooting, singular endpoints, inconsistent boundary counts, poor initial guesses, mesh underresolution, ill-conditioned jacobians, and arbitrary normalization of eigenfunctions.  Use scaled tolerances and report both success criteria and diagnostic evidence.}
\noterow{Implementation checklist}{\begin{itemize}
\item Validate dimensions, domains, ordering, and structural assumptions.
\item Guard small divisors, invalid states, overflow, underflow, and nonfinite values.
\item Make mutation, workspace, indexing, and termination behavior explicit.
\item Test a simple exact case, a difficult valid case, a boundary case, and an expected failure.
\end{itemize}}
\end{cornell}
\begin{summarybox}{Section summary}
A boundary-value problem imposes conditions at separated points, so the entire solution trajectory must satisfy coupled endpoint constraints.  Select this technique only when its assumptions match the data, and accept a result only after an independent residual, error estimate, invariant, or refinement check.
\end{summarybox}

\section{18.1 The Shooting Method}
\begin{cornell}
\noterow{What is the central idea?}{Shooting converts unknown initial conditions into parameters and solves an endpoint residual equation after each initial-value integration.}
\noterow{How should it be used?}{For The Shooting Method, begin from explicit inputs, outputs, preconditions, and representation choices.  Count differential order and independent boundary constraints; choose shooting, multiple shooting, matching, or relaxation; provide initial guesses; solve endpoint or mesh residuals; refine and verify conditions.}
\noterow{Quantitative anchor}{\[
R(s)=B\bigl(y(b;s)\bigr)=0
\]
State the dimensions, scaling, and validity assumptions before using this relation.  Check the resulting residual or invariant rather than trusting an iteration count alone.}
\noterow{Cost and reliability}{Analyze arithmetic work, storage, data movement, and convergence separately.  Relevant failure pressures include exponentially unstable shooting, singular endpoints, inconsistent boundary counts, poor initial guesses, mesh underresolution, ill-conditioned jacobians, and arbitrary normalization of eigenfunctions.  Use scaled tolerances and report both success criteria and diagnostic evidence.}
\noterow{Implementation checklist}{\begin{itemize}
\item Validate dimensions, domains, ordering, and structural assumptions.
\item Guard small divisors, invalid states, overflow, underflow, and nonfinite values.
\item Make mutation, workspace, indexing, and termination behavior explicit.
\item Test a simple exact case, a difficult valid case, a boundary case, and an expected failure.
\end{itemize}}
\end{cornell}
\begin{summarybox}{Section summary}
Shooting converts unknown initial conditions into parameters and solves an endpoint residual equation after each initial-value integration.  Select this technique only when its assumptions match the data, and accept a result only after an independent residual, error estimate, invariant, or refinement check.
\end{summarybox}

\section{18.2 Shooting to a Fitting Point}
\begin{cornell}
\noterow{What is the central idea?}{Integrating from both boundaries toward an interior match point can reduce instability that grows across the full interval.}
\noterow{How should it be used?}{For Shooting to a Fitting Point, begin from explicit inputs, outputs, preconditions, and representation choices.  Count differential order and independent boundary constraints; choose shooting, multiple shooting, matching, or relaxation; provide initial guesses; solve endpoint or mesh residuals; refine and verify conditions.}
\noterow{Quantitative anchor}{\[
R(s)=B\bigl(y(b;s)\bigr)=0
\]
State the dimensions, scaling, and validity assumptions before using this relation.  Check the resulting residual or invariant rather than trusting an iteration count alone.}
\noterow{Cost and reliability}{Analyze arithmetic work, storage, data movement, and convergence separately.  Relevant failure pressures include exponentially unstable shooting, singular endpoints, inconsistent boundary counts, poor initial guesses, mesh underresolution, ill-conditioned jacobians, and arbitrary normalization of eigenfunctions.  Use scaled tolerances and report both success criteria and diagnostic evidence.}
\noterow{Implementation checklist}{\begin{itemize}
\item Validate dimensions, domains, ordering, and structural assumptions.
\item Guard small divisors, invalid states, overflow, underflow, and nonfinite values.
\item Make mutation, workspace, indexing, and termination behavior explicit.
\item Test a simple exact case, a difficult valid case, a boundary case, and an expected failure.
\end{itemize}}
\end{cornell}
\begin{summarybox}{Section summary}
Integrating from both boundaries toward an interior match point can reduce instability that grows across the full interval.  Select this technique only when its assumptions match the data, and accept a result only after an independent residual, error estimate, invariant, or refinement check.
\end{summarybox}

\section{18.3 Relaxation Methods}
\begin{cornell}
\noterow{What is the central idea?}{Relaxation discretizes the entire trajectory and solves the resulting coupled nonlinear algebraic equations simultaneously.}
\noterow{How should it be used?}{For Relaxation Methods, begin from explicit inputs, outputs, preconditions, and representation choices.  Count differential order and independent boundary constraints; choose shooting, multiple shooting, matching, or relaxation; provide initial guesses; solve endpoint or mesh residuals; refine and verify conditions.}
\noterow{Quantitative lens}{Identify the governing equation, recurrence, probability law, or geometric predicate for Relaxation Methods.  Define every variable and scale; then derive a residual, error estimate, or invariant that can be checked independently.}
\noterow{Cost and reliability}{Analyze arithmetic work, storage, data movement, and convergence separately.  Relevant failure pressures include exponentially unstable shooting, singular endpoints, inconsistent boundary counts, poor initial guesses, mesh underresolution, ill-conditioned jacobians, and arbitrary normalization of eigenfunctions.  Use scaled tolerances and report both success criteria and diagnostic evidence.}
\noterow{Implementation checklist}{\begin{itemize}
\item Validate dimensions, domains, ordering, and structural assumptions.
\item Guard small divisors, invalid states, overflow, underflow, and nonfinite values.
\item Make mutation, workspace, indexing, and termination behavior explicit.
\item Test a simple exact case, a difficult valid case, a boundary case, and an expected failure.
\end{itemize}}
\end{cornell}
\begin{summarybox}{Section summary}
Relaxation discretizes the entire trajectory and solves the resulting coupled nonlinear algebraic equations simultaneously.  Select this technique only when its assumptions match the data, and accept a result only after an independent residual, error estimate, invariant, or refinement check.
\end{summarybox}

\section{18.4 A Worked Example: Spheroidal Harmonics}
\begin{cornell}
\noterow{What is the central idea?}{The example integrates eigenvalue adjustment, endpoint behavior, normalization, and matching into a complete boundary-value workflow.}
\noterow{How should it be used?}{For A Worked Example: Spheroidal Harmonics, begin from explicit inputs, outputs, preconditions, and representation choices.  Count differential order and independent boundary constraints; choose shooting, multiple shooting, matching, or relaxation; provide initial guesses; solve endpoint or mesh residuals; refine and verify conditions.}
\noterow{Quantitative lens}{Identify the governing equation, recurrence, probability law, or geometric predicate for A Worked Example: Spheroidal Harmonics.  Define every variable and scale; then derive a residual, error estimate, or invariant that can be checked independently.}
\noterow{Cost and reliability}{Analyze arithmetic work, storage, data movement, and convergence separately.  Relevant failure pressures include exponentially unstable shooting, singular endpoints, inconsistent boundary counts, poor initial guesses, mesh underresolution, ill-conditioned jacobians, and arbitrary normalization of eigenfunctions.  Use scaled tolerances and report both success criteria and diagnostic evidence.}
\noterow{Implementation checklist}{\begin{itemize}
\item Validate dimensions, domains, ordering, and structural assumptions.
\item Guard small divisors, invalid states, overflow, underflow, and nonfinite values.
\item Make mutation, workspace, indexing, and termination behavior explicit.
\item Test a simple exact case, a difficult valid case, a boundary case, and an expected failure.
\end{itemize}}
\end{cornell}
\begin{summarybox}{Section summary}
The example integrates eigenvalue adjustment, endpoint behavior, normalization, and matching into a complete boundary-value workflow.  Select this technique only when its assumptions match the data, and accept a result only after an independent residual, error estimate, invariant, or refinement check.
\end{summarybox}

\section{18.5 Automated Allocation of Mesh Points}
\begin{cornell}
\noterow{What is the central idea?}{Adaptive meshes concentrate nodes where the solution changes rapidly while maintaining a smooth mapping and a meaningful error indicator.}
\noterow{How should it be used?}{For Automated Allocation of Mesh Points, begin from explicit inputs, outputs, preconditions, and representation choices.  Count differential order and independent boundary constraints; choose shooting, multiple shooting, matching, or relaxation; provide initial guesses; solve endpoint or mesh residuals; refine and verify conditions.}
\noterow{Quantitative lens}{Identify the governing equation, recurrence, probability law, or geometric predicate for Automated Allocation of Mesh Points.  Define every variable and scale; then derive a residual, error estimate, or invariant that can be checked independently.}
\noterow{Cost and reliability}{Analyze arithmetic work, storage, data movement, and convergence separately.  Relevant failure pressures include exponentially unstable shooting, singular endpoints, inconsistent boundary counts, poor initial guesses, mesh underresolution, ill-conditioned jacobians, and arbitrary normalization of eigenfunctions.  Use scaled tolerances and report both success criteria and diagnostic evidence.}
\noterow{Implementation checklist}{\begin{itemize}
\item Validate dimensions, domains, ordering, and structural assumptions.
\item Guard small divisors, invalid states, overflow, underflow, and nonfinite values.
\item Make mutation, workspace, indexing, and termination behavior explicit.
\item Test a simple exact case, a difficult valid case, a boundary case, and an expected failure.
\end{itemize}}
\end{cornell}
\begin{summarybox}{Section summary}
Adaptive meshes concentrate nodes where the solution changes rapidly while maintaining a smooth mapping and a meaningful error indicator.  Select this technique only when its assumptions match the data, and accept a result only after an independent residual, error estimate, invariant, or refinement check.
\end{summarybox}

\section{18.6 Handling Internal Boundary Conditions or Singular Points}
\begin{cornell}
\noterow{What is the central idea?}{Interior constraints and singular points require domain splitting, matching conditions, or analytic asymptotic information.}
\noterow{How should it be used?}{For Handling Internal Boundary Conditions or Singular Points, begin from explicit inputs, outputs, preconditions, and representation choices.  Count differential order and independent boundary constraints; choose shooting, multiple shooting, matching, or relaxation; provide initial guesses; solve endpoint or mesh residuals; refine and verify conditions.}
\noterow{Quantitative lens}{Identify the governing equation, recurrence, probability law, or geometric predicate for Handling Internal Boundary Conditions or Singular Points.  Define every variable and scale; then derive a residual, error estimate, or invariant that can be checked independently.}
\noterow{Cost and reliability}{Analyze arithmetic work, storage, data movement, and convergence separately.  Relevant failure pressures include exponentially unstable shooting, singular endpoints, inconsistent boundary counts, poor initial guesses, mesh underresolution, ill-conditioned jacobians, and arbitrary normalization of eigenfunctions.  Use scaled tolerances and report both success criteria and diagnostic evidence.}
\noterow{Implementation checklist}{\begin{itemize}
\item Validate dimensions, domains, ordering, and structural assumptions.
\item Guard small divisors, invalid states, overflow, underflow, and nonfinite values.
\item Make mutation, workspace, indexing, and termination behavior explicit.
\item Test a simple exact case, a difficult valid case, a boundary case, and an expected failure.
\end{itemize}}
\end{cornell}
\begin{summarybox}{Section summary}
Interior constraints and singular points require domain splitting, matching conditions, or analytic asymptotic information.  Select this technique only when its assumptions match the data, and accept a result only after an independent residual, error estimate, invariant, or refinement check.
\end{summarybox}

\section*{Method comparison}
\begin{center}
\small
\begin{tabularx}{\linewidth}{>{\bfseries\raggedright\arraybackslash}p{0.22\linewidth}XX}
\toprule
Method or lens & Best use & Main caution \\
\midrule
Single shooting & Stable IVP and few unknown initial values & Simple; can amplify instability \\
Multiple shooting & Long or unstable intervals & More variables; better conditioning \\
Relaxation & Strongly coupled or singular BVPs & Large sparse nonlinear system \\
\bottomrule
\end{tabularx}
\end{center}

\section*{Worked example}
\begin{exambox}{Independent worked example}
Solve $y''=0$ with $y(0)=1$ and $y(2)=5$.  Writing $y=a+bx$ gives $a=1$ and $b=2$, so $y=1+2x$.  Both boundary residuals vanish.  A shooting implementation recovers the unknown initial slope $2$; for unstable equations, small slope error can grow dramatically.
\end{exambox}

\section*{Chapter synthesis}
\begin{summarybox}{Chapter synthesis}
The chapter's methods should be viewed as a decision system rather than a list of routines.  Begin with the mathematical problem and its conditioning, exploit structure, select an algorithm with compatible stability and cost, and verify the computed answer with evidence that is independent of the internal iteration.  The recurring implementation discipline is: count differential order and independent boundary constraints; choose shooting, multiple shooting, matching, or relaxation; provide initial guesses; solve endpoint or mesh residuals; refine and verify conditions.
\end{summarybox}

\subsection*{Common mistakes and numerical hazards}
\begin{warnbox}{Common mistakes and numerical hazards}
Exponentially unstable shooting, singular endpoints, inconsistent boundary counts, poor initial guesses, mesh underresolution, ill-conditioned Jacobians, and arbitrary normalization of eigenfunctions.  A second recurring mistake is reporting only a computed value: always retain tolerances, iteration counts, residuals, refinement behavior, and relevant structural checks.
\end{warnbox}

\subsection*{Retrieval practice}
\textbf{Questions}
\begin{enumerate}
\item What problem does The Shooting Method address, and which assumption is easiest to violate?
\item What problem does Shooting to a Fitting Point address, and which assumption is easiest to violate?
\item What problem does Relaxation Methods address, and which assumption is easiest to violate?
\item What problem does Automated Allocation of Mesh Points address, and which assumption is easiest to violate?
\item What problem does Handling Internal Boundary Conditions or Singular Points address, and which assumption is easiest to violate?
\item How do conditioning, stability, and the final residual play different roles in this chapter?
\end{enumerate}

\textbf{Brief answers}
\begin{enumerate}
\item Shooting converts unknown initial conditions into parameters and solves an endpoint residual equation after each initial-value integration. Recheck the section's domain, structure, scaling, and failure diagnostics.
\item Integrating from both boundaries toward an interior match point can reduce instability that grows across the full interval. Recheck the section's domain, structure, scaling, and failure diagnostics.
\item Relaxation discretizes the entire trajectory and solves the resulting coupled nonlinear algebraic equations simultaneously. Recheck the section's domain, structure, scaling, and failure diagnostics.
\item Adaptive meshes concentrate nodes where the solution changes rapidly while maintaining a smooth mapping and a meaningful error indicator. Recheck the section's domain, structure, scaling, and failure diagnostics.
\item Interior constraints and singular points require domain splitting, matching conditions, or analytic asymptotic information. Recheck the section's domain, structure, scaling, and failure diagnostics.
\item Conditioning describes sensitivity of the mathematical problem, stability describes error propagation by the algorithm, and the residual measures satisfaction of the computed equations; none is a universal substitute for the others.
\end{enumerate}

\subsection*{Mastery checklist}
\begin{itemize}
\item[$\square$] I can explain and select The Shooting Method without relying on a memorized code listing.
\item[$\square$] I can explain and select Shooting to a Fitting Point without relying on a memorized code listing.
\item[$\square$] I can explain and select Relaxation Methods without relying on a memorized code listing.
\item[$\square$] I can explain and select A Worked Example: Spheroidal Harmonics without relying on a memorized code listing.
\item[$\square$] I can explain and select Automated Allocation of Mesh Points without relying on a memorized code listing.
\item[$\square$] I can state a scaled stopping rule and an independent verification check.
\item[$\square$] I can identify at least one conditioning risk and one algorithmic stability risk.
\item[$\square$] I can estimate time, storage, and data-movement costs at the appropriate level.
\end{itemize}

\end{document}

